\section{Evaluation Protocol}
\label{sec:evaluation}

Two different notions of ``holdout'' are used, and keeping them distinct is
the whole point of this report.

\subsection{Warm-start: leave-last-one-out}

Within the population of users retained after core filtering, each user's
last interaction (by true review timestamp) is held out for test, the
second-to-last for validation, and the rest for training. Users without
enough history to leave at least three training interactions contribute
only to training and are never evaluated, which avoids inflating metrics
with users whose ``history'' is a single review.

\subsection{Cold-start: entities withheld entirely}

Before the warm-start split is computed, a fraction of users and a fraction
of items (8\% each, by default) are selected uniformly at random and every
interaction touching either is removed from the training-eligible pool
\emph{before} any splitting happens. The warm-start split above is then
computed only over the remainder. Three evaluation slices follow directly
from this construction:

\begin{itemize}
  \item \textbf{New user, known item} (\texttt{cold\_user\_eval}): reviews
        by a withheld user against a product that \emph{was} trained on.
        Tests whether shared demographic features let the model retrieve
        for a user it has never seen.
  \item \textbf{Known user, new item} (\texttt{cold\_item\_eval}): reviews
        by a trained user against a withheld product. Tests whether shared
        brand/category/ingredient features let the model retrieve a product
        it has never seen for a user whose taste it has learned.
  \item \textbf{New user, new item} (\texttt{cold\_both\_eval}): both sides
        withheld -- the hardest slice, reported for completeness.
\end{itemize}

Because the identity-feature vocabulary (Section~\ref{sec:features}) is
built from the training partition alone, a withheld entity provably has no
identity token in either the hybrid model or the ID-only baseline. The
exclusion set passed to every model's ranking call is empty for these
slices: there is no training history to hide, which is exactly the point.

\subsection{Metrics}

For every evaluated user, the model ranks the full catalogue (excluding
training history, where applicable) and NDCG@$K$, Recall@$K$ and MAP@$K$ are
computed against the relevant set -- a single item for the warm-start
partitions (leave-last-one-out), potentially several for the cold-start
slices (every interaction a withheld entity has is relevant, since all of
them were excluded from training).
\begin{align}
  \text{NDCG@}K &= \frac{1}{|\Uset_{\text{eval}}|} \sum_{u}
    \frac{\sum_{k=1}^{K} \mathbb{1}[\text{rel}(k)] / \log_2(k+1)}{\text{IDCG}_u} , \\
  \text{Recall@}K &= \frac{1}{|\Uset_{\text{eval}}|} \sum_{u}
    \frac{|\{\text{relevant items retrieved in top-}K\}|}{|\text{relevant items for } u|} , \\
  \text{MAP@}K &= \frac{1}{|\Uset_{\text{eval}}|} \sum_{u}
    \frac{1}{\min(|\text{relevant}_u|, K)} \sum_{k=1}^{K} \mathbb{1}[\text{rel}(k)] \cdot \text{Precision@}k .
\end{align}

\subsection{Hyperparameter selection}

Hyperparameters for every model (hybrid and baselines alike) are fixed
before the test partition and both cold-start slices are read; the values
used are recorded in \texttt{src/recommender/cli.py}
(\texttt{BEST\_HYBRID}, \texttt{BEST\_ID\_MF}) and are the same values the
table-generation script (\texttt{scripts/export\_results\_table.py}) uses,
so the numbers in Section~\ref{sec:results} are reproducible from a clean
checkout with \texttt{python scripts/download\_data.py \&\& python
scripts/export\_results\_table.py}.
